
\documentclass[12pt,a4paper]{article}

% ======================================
% PACOTES
% ======================================
\usepackage[utf8]{inputenc}
\usepackage[T1]{fontenc}
\usepackage[brazil]{babel}
\usepackage{geometry}
\usepackage{graphicx}
\usepackage{float}
\usepackage{booktabs}
\usepackage{setspace}
\usepackage{indentfirst}
\usepackage{hyperref}

% ======================================
% FORMATAÇÃO
% ======================================
\geometry{
    left=3cm,
    right=2cm,
    top=3cm,
    bottom=2cm
}

\onehalfspacing
\setlength{\parindent}{1.25cm}
\setlength{\parskip}{0pt}

\hypersetup{
    colorlinks=true,
    linkcolor=black,
    citecolor=black,
    urlcolor=blue
}

\begin{document}

% ======================================
% CAPA
% ======================================
\begin{titlepage}
    \begin{center}

        \textbf{INSTITUTO FEDERAL CATARINENSE}\\
        \textbf{CAMPUS RIO DO SUL}

        \vspace{2cm}

        \textbf{BACHARELADO EM CIÊNCIA DA COMPUTAÇÃO}

        \vspace{3cm}

        {\Large\bfseries
        LABORATÓRIO DE COMUNICAÇÃO\\
        PUBLISH/SUBSCRIBE COM AMQP\\
        E ACTIVEMQ ARTEMIS
        }

        \vspace{3cm}

        \textbf{Acadêmico:} Nome do acadêmico\\
        \textbf{Professor:} Nome do professor\\
        \textbf{Disciplina:} Desenvolvimento Web

        \vfill

        Rio do Sul -- SC\\
        2026

    \end{center}
\end{titlepage}

% ======================================
% RESUMO
% ======================================
\section*{Resumo}
\addcontentsline{toc}{section}{Resumo}

Este relatório apresenta a implementação e a análise
de experimentos de comunicação distribuída utilizando
o protocolo AMQP 1.0, o broker Apache ActiveMQ Artemis
e aplicações desenvolvidas em Java.

Foram realizados experimentos envolvendo a distribuição
de mensagens em filas, a comunicação Publish/Subscribe
com múltiplos assinantes e o armazenamento persistente
de mensagens.

Os resultados permitem observar diferentes estratégias
de distribuição e entrega de mensagens em sistemas
distribuídos.

\textbf{Palavras-chave:} AMQP; ActiveMQ Artemis;
Publish/Subscribe; Java; Sistemas Distribuídos.

\newpage

\tableofcontents
\newpage

% ======================================
% 1 INTRODUÇÃO
% ======================================
\section{Introdução}

A comunicação entre aplicações distribuídas
pode ser realizada por meio de sistemas de
mensageria, nos quais produtores e consumidores
trocam mensagens utilizando um intermediário
denominado \textit{broker}.

Neste laboratório, foi utilizado o Apache
ActiveMQ Artemis para implementar experimentos
de comunicação utilizando o protocolo AMQP 1.0.

O objetivo foi investigar o comportamento
de filas, tópicos e mensagens persistentes
por meio de aplicações Java executadas
simultaneamente.

\subsection{Objetivos}

O objetivo geral consiste em implementar e
avaliar diferentes mecanismos de comunicação
por mensagens utilizando AMQP e ActiveMQ.

Os objetivos específicos são:

\begin{itemize}
    \item Configurar um broker ActiveMQ em Docker;
    \item Implementar produtores e consumidores em Java;
    \item Avaliar a distribuição de mensagens em filas;
    \item Demonstrar a comunicação fan-out;
    \item Verificar a persistência após reinício do broker.
\end{itemize}

% ======================================
% 2 REFERENCIAL TEÓRICO
% ======================================
\section{Referencial Teórico}

\subsection{Sistemas de Mensageria}

% Desenvolver com base nos artigos indicados.

\subsection{Arquitetura Publish/Subscribe}

% Explicar produtores, broker e assinantes.

\subsection{Protocolo AMQP}

% Explicar AMQP 1.0.

\subsection{Apache ActiveMQ Artemis}

% Explicar o funcionamento do broker.

\subsection{Filas, Tópicos e Persistência}

% Explicar ANYCAST, MULTICAST e persistência.

% ======================================
% 3 METODOLOGIA
% ======================================
\section{Metodologia}

\subsection{Ambiente Experimental}

Os experimentos foram realizados em um ambiente
Linux, utilizando as seguintes tecnologias:

\begin{itemize}
    \item Sistema operacional: Arch Linux;
    \item Linguagem de programação: Java;
    \item Gerenciador de dependências: Apache Maven;
    \item Broker: Apache ActiveMQ Artemis;
    \item Protocolo: AMQP 1.0;
    \item Biblioteca: Apache Qpid JMS;
    \item Contêineres: Docker;
    \item IDE: IntelliJ IDEA.
\end{itemize}

O broker foi executado em um contêiner Docker,
disponibilizando o serviço de mensageria
na porta 61616.

As aplicações Java estabeleceram conexões
com o broker por meio da biblioteca Qpid JMS.

\subsection{Organização dos Experimentos}

O laboratório foi dividido em três experimentos:
distribuição de mensagens em fila, comunicação
fan-out por tópicos e persistência de mensagens.

% ======================================
% 4 EXPERIMENTO I
% ======================================
\section{Experimento I: Filas e Distribuição de Mensagens}

\subsection{Implementação}

Foi implementado um produtor Java responsável
por enviar cinco mensagens para uma fila
gerenciada pelo ActiveMQ Artemis.

As mensagens foram publicadas com intervalos
de 60 segundos.

Também foram implementados consumidores
responsáveis por receber individualmente
as mensagens disponibilizadas pelo broker.

\subsection{Resultados}

Na comunicação baseada em filas com roteamento
ANYCAST, uma mensagem é direcionada a um dos
consumidores associados à fila.

O experimento demonstrou a recepção de mensagens
por diferentes processos Java, identificados
por seus respectivos PIDs.

% Inserir prints dos consumidores e da fila.

\subsection{Discussão}

O mecanismo de filas permite distribuir
mensagens entre consumidores, diferentemente
do modelo fan-out, no qual uma publicação
pode alcançar múltiplos assinantes.

% ======================================
% 5 EXPERIMENTO II
% ======================================
\section{Experimento II: Fan-out com Tópicos}

\subsection{Implementação}

Foi utilizado um tópico para implementar
a comunicação Publish/Subscribe.

Cinco processos assinantes foram conectados
simultaneamente ao broker.

Em seguida, um produtor publicou cinco
mensagens, utilizando intervalos de
60 segundos entre as publicações.

\subsection{Resultados}

Todos os cinco assinantes receberam as
cinco mensagens publicadas, na ordem
correspondente aos eventos numerados
de 1 até 5.

O experimento apresentou um total de
25 recebimentos.

\begin{table}[H]
\centering
\caption{Resultados do experimento fan-out}
\begin{tabular}{lc}
\toprule
\textbf{Indicador} & \textbf{Resultado} \\
\midrule
Mensagens publicadas & 5 \\
Assinantes simultâneos & 5 \\
Intervalo entre publicações & 60 s \\
Mensagens por assinante & 5 \\
Total de recebimentos & 25 \\
\bottomrule
\end{tabular}
\label{tab:fanout}
\end{table}

\subsection{Discussão}

Os resultados demonstram o comportamento
MULTICAST da comunicação Publish/Subscribe,
em que cada assinante ativo recebe uma
cópia das mensagens publicadas no tópico.

% Inserir evidências e aprofundar a análise.

% ======================================
% 6 EXPERIMENTO III
% ======================================
\section{Experimento III: Persistência de Mensagens}

\subsection{Implementação}

Foi implementado um produtor utilizando
o modo de entrega \texttt{PERSISTENT}
da API Jakarta Messaging.

A mensagem foi enviada para a fila
\texttt{fila.persistente}.

Após o envio, o contêiner responsável
pelo ActiveMQ Artemis foi reiniciado.

\subsection{Resultados}

A mensagem permaneceu armazenada na fila
após o reinício do broker, demonstrando
a preservação da mensagem persistente
nesse cenário experimental.

% Inserir print antes e depois do reinício.
% Documentar o consumo posterior da mensagem.

\subsection{Discussão}

A persistência permite que mensagens
não consumidas sejam preservadas pelo broker
durante determinadas situações de reinício
ou interrupção do serviço.

Esse mecanismo é relevante para aplicações
que exigem maior confiabilidade na
entrega de mensagens.

% ======================================
% 7 ANÁLISE COMPARATIVA
% ======================================
\section{Análise Comparativa}

\begin{table}[H]
\centering
\caption{Comparação dos mecanismos experimentados}
\begin{tabular}{lll}
\toprule
\textbf{Mecanismo} &
\textbf{Destino} &
\textbf{Comportamento} \\
\midrule
ANYCAST & Fila & Distribuição \\
MULTICAST & Tópico & Fan-out \\
PERSISTENT & Fila & Armazenamento \\
\bottomrule
\end{tabular}
\label{tab:comparacao}
\end{table}

Os mecanismos avaliados apresentam
características distintas de distribuição
e armazenamento de mensagens.

Enquanto as filas permitem compartilhar
a carga de processamento entre consumidores,
os tópicos possibilitam distribuir uma
mesma publicação a múltiplos assinantes.

Já a persistência está relacionada
à preservação das mensagens armazenadas,
não sendo uma alternativa aos mecanismos
de roteamento.

% ======================================
% 8 CONCLUSÃO
% ======================================
\section{Conclusão}

Os experimentos permitiram observar
diferentes comportamentos da comunicação
por mensagens utilizando AMQP 1.0
e Apache ActiveMQ Artemis.

O experimento com filas demonstrou a
distribuição de mensagens entre consumidores.

No experimento fan-out, cinco assinantes
receberam cinco publicações cada,
totalizando 25 recebimentos.

O teste de persistência demonstrou que
a mensagem armazenada permaneceu disponível
após o reinício do broker.

Os resultados evidenciam a importância
da escolha do mecanismo de comunicação
conforme os requisitos das aplicações
distribuídas.

% ======================================
% REFERÊNCIAS
% ======================================
\newpage

\bibliographystyle{plain}
\bibliography{referencias}

\end{document}
